\section{The optimizations, one by one}
\label{sec:optimizations}

The first working version already decoded image \texttt{a002} seven times faster than NBIS,
because it vectorized the transform from the start. Table~\ref{tab:history} follows the decoder
from there; each row was measured right after the change, on the same image and core, so the
numbers carry the noise of separate runs (a few percent). Two changes that did not pay are
listed as well.

\begin{table}[ht]
  \centering\small
  \begin{tabular}{@{}rp{9.6cm}rr@{}}
    \toprule
    & change & AVX2 (ms) & vs NBIS \\
    \midrule
    0 & NBIS 5.0.0, for reference & 53.5 & 1.0$\times$ \\
    1 & first version: NBIS schedules evaluated across lines, transposes for the row pass & 7.59 & 7.0$\times$ \\
    2 & schedules cached per image size; buffers reused; pixel conversion vectorized & 6.17 & 8.7$\times$ \\
    3 & entropy decoding writes straight into the coefficient array & 5.78 & 9.3$\times$ \\
    -- & eight accumulators instead of four in the column kernel (no effect, reverted) & 5.86 & \\
    4 & row-direct synthesis, two dependency chains (no effect) & 6.07 & \\
    5 & row-direct synthesis, eight dependency chains & 5.16 & 10.4$\times$ \\
    -- & template column kernel (slower with 8 lanes; kept for 4-lane backends only) & 5.45 & \\
    6 & column kernel chosen by vector width & 4.55 & 11.8$\times$ \\
    7 & entropy decoding into a linear buffer without branches, vector dequantization & 4.48 & 11.9$\times$ \\
    8 & four threads: three entropy-coded blocks and the large splits in parallel & 2.47 & 21.7$\times$ \\
    \bottomrule
  \end{tabular}
  \caption{Decoding \texttt{a002} at 0.75 bits per pixel, exact mode, AVX2, as the decoder evolved.}
  \label{tab:history}
\end{table}

\subsection{Fixed costs}

The first profile (stage times of version 1) showed 0.64~ms spent building the transform
schedules, 1.19~ms converting coefficients to pixels with scalar code and a fresh 5.8~MB
coefficient array per image, whose page faults spread over the other stages. Caching the
schedules by line length and filter bank, keeping the buffers in the decoder object, and a
vector pixel kernel brought the first 1.4~ms down.

\subsection{Entropy decoding}

Huffman decoding is sequential within a block: each code's length is known only after it is
read. A bare loop that decodes symbols and discards them runs at 3.3~ns per symbol, but placing
the symbols cost another 4.4~ns: the symbol type (a run of zeros or a value) is unpredictable,
and every value write needed the bookkeeping of which subband row it belonged to. Two changes
fixed this.
\begin{enumerate}
  \item A 10-bit lookup table whose entries hold the code length, the symbol's kind and its
    value, so that 98\% of the symbols cost one table load (codes longer than 10 bits, about
    2\%, take a slower path).
  \item A first pass that writes indices, in stream order, into a linear buffer with one code
    path for runs and values (Algorithm~\ref{alg:entropy}), and a second, vectorized pass that
    dequantizes each subband row into its place in the coefficient array, without branches.
\end{enumerate}
The entropy stage of \texttt{a002} went from 1.62~ms to 0.63~ms (2.6$\times$); at 2.25 bits per
pixel, with four times the symbols, from 5.78~ms to 2.40~ms.

\begin{algorithm}[ht]
  \caption{Branch-free placement of decoded symbols.}
  \label{alg:entropy}
  \begin{algorithmic}[1]
    \Require a block's entropy-coded bits; a zeroed buffer $L$ with one spare slot; range $[k_0, k_1)$
    \State $k \gets k_0$
    \While{$k < k_1$}
      \State $(\ell, \mathit{kind}, v) \gets \mathrm{LUT}[\text{next 10 bits}]$;\ consume $\ell$ bits \Comment{rare: long codes, escapes}
      \State $\mathit{run} \gets v$ \textbf{if} $\mathit{kind} = \text{run}$ \textbf{else} $0$;\quad $\mathit{val} \gets v$ \textbf{if} $\mathit{kind} = \text{value}$ \textbf{else} $0$
      \State $k \gets k + \mathit{run}$
      \State $L[k] \gets \mathit{val}$ \Comment{a run writes a zero where the next value goes}
      \State $k \gets k + [\mathit{val} \ne 0]$
    \EndWhile
  \end{algorithmic}
\end{algorithm}

\subsection{The inverse transform}

After the fixed costs, the transform took two thirds of the time, and the vector version was
only twice as fast as the scalar one with eight lanes. Two experiments located the problem.
Replacing every schedule by the identity, which leaves only the data movement, cost 2.5 of the
4.0~ms: the transform is limited by moving data, not by arithmetic. And the fast mode, which
does half the multiplications, was not faster. With the identity schedules, the row pass,
which transposed blocks of eight rows into a scratch buffer and back, cost twice as much as the
column pass.

The row-direct kernel (Section~\ref{sec:exactness}) removed the transposes of the synthesis but
was not faster at first: each output accumulates about eight products, every addition waits
four cycles for the previous one, and only two chains (even and odd) were in flight. Evaluating
four vectors of each parity at once gives eight independent chains and brought the transform of
\texttt{a002} from 3.99 to 3.13~ms. The same trick made the column pass slower on AVX2, where
the strip of 64 columns already provides enough independent work, so each backend uses the
kernel that suits its width. Fusing the column and the row pass of a split, so that the
intermediate stays in cache, was considered and rejected: synthesis cannot run in place (the
lowpass and highpass halves are read far from where the outputs go), and the copy it would need
moves as much data as it saves.

\subsection{Threads}

Inside one image, the three entropy-coded blocks of encoder number two are independent, and
their positions in the linear buffer are known from the subband sizes, so three threads decode
them. The transform's large splits (the image and its quadrants) share columns and then rows
among threads, separated by a barrier; the small, deep splits run on one thread, where
synchronization would cost more than the work.

\subsection{The encoder}

The same ideas, plus four specific to encoding: quantization divides once per coefficient with
the truncation identity of Section~\ref{sec:exactness} instead of branching per sign (0.72 to
0.27~ms on \texttt{cmp00001}); normalization converts bytes to floats eight at a time; the
tokenizer skips zeros a vector at a time, testing eight or sixteen coefficients with one
comparison; and the bit writer emits code and extra bits in one call and four bytes at a time.
The encoder of \texttt{a002} went from 10.3~ms in its first version to 6.9~ms.
